% Alur penelitian tugas akhir (TikZ).
% Di-input dari BAB III §3.2 Alur Tugas Akhir.

\begin{figure}[H]
\centering
\resizebox{\textwidth}{!}{%
\begin{tikzpicture}[
  proc/.style={draw=black, fill=white, rectangle, rounded corners=3pt,
    text width=4.25cm, minimum height=1.15cm, inner sep=4pt,
    align=center, font=\footnotesize},
  phase/.style={draw=gray!70, fill=gray!3, rectangle, rounded corners=10pt,
    minimum width=5.2cm, minimum height=5.4cm},
  startstop/.style={draw=black, fill=white, rectangle, rounded corners=8pt,
    minimum width=2.2cm, minimum height=0.8cm, align=center, font=\footnotesize},
  ar/.style={-{Stealth}, thick},
  feedback/.style={-{Stealth}, thick, dashed},
  lab/.style={font=\scriptsize, midway, align=center},
]

%--- Kontainer fase (digambar lebih dahulu agar berada di belakang node) ---
\node[phase] (prep) at (0, 2.9) {};
\node[phase] (dev)  at (6, 2.9) {};
\node[phase] (test) at (12, 2.9) {};

\node[font=\scriptsize, anchor=north west] at ([xshift=0.2cm,yshift=-0.18cm]prep.north west)
  {Persiapan};
\node[font=\scriptsize, anchor=north west] at ([xshift=0.2cm,yshift=-0.18cm]dev.north west)
  {Pengembangan};
\node[font=\scriptsize, anchor=north west] at ([xshift=0.2cm,yshift=-0.18cm]test.north west)
  {Pengujian};

%--- Persiapan ---
\node[startstop] (start) at (0, 6.5) {Mulai};
\node[proc] (s1) at (0, 4.2)
  {Studi Literatur \&\\Analisis Dokumen Bank Indonesia};
\node[proc] (s2) at (0, 2.1)
  {Perancangan Sistem\\(\textit{Activity Diagram}, Arsitektur, \textit{Use Case})};

%--- Pengembangan ---
\node[proc, minimum height=1.55cm] (s3) at (6, 3.1)
  {Implementasi Sistem\\(\textit{Chaincode}, Jaringan Fabric, \textit{Backend} API)};

%--- Pengujian ---
\node[proc] (s4) at (12, 4.2)
  {Pengujian\\(Fungsional + Kinerja dengan Hyperledger Caliper)};
\node[proc] (s5) at (12, 2.1)
  {Analisis Hasil \&\\Penarikan Kesimpulan};
\node[startstop] (finish) at (12, -0.6) {Selesai};

%--- Alur utama ---
\draw[ar] (start.south) -- (s1.north);
\draw[ar] (s1.south) -- (s2.north);
\draw[ar] (s2.east) -- (s3.west);
\draw[ar] (s3.east) -- (s4.west);
\draw[ar] (s4.south) -- (s5.north);
\draw[ar] (s5.south) -- (finish.north);

%--- Umpan balik iteratif (prototype loop) ---
\draw[feedback] (s4.north west) to[out=145,in=35,looseness=1.35]
  node[lab,below,pos=0.5]{revisi jika perlu} (s3.north east);

\end{tikzpicture}%
}
\caption{Alur pelaksanaan tugas akhir: dari studi literatur hingga penarikan kesimpulan, dengan umpan balik iteratif antara implementasi dan pengujian}
\label{fig:alur-penelitian}
\end{figure}
